\section*{Question 4}
Two charges are placed in two different arrangements as shown below: 

\tdplotsetmaincoords{70}{110}
\begin{figure}[h!]
\centering
\begin{tikzpicture}[tdplot_main_coords, scale=1.4]

% Axes
\draw[->] (0,0,0) -- (0,0,2) node[anchor=south]{$z$};
\draw[->] (0,0,0) -- (2,0,0) node[anchor=north east]{$x$};
\draw[->] (0,0,0) -- (0,2,0) node[anchor=west]{$y$};

% Arrangement (a)
% Labels
\node at (1,-1,2.) {\textbf{(a)}};


% Charge at (0,0,a)
\filldraw[red] (0,0,1.5) circle (2pt) node[anchor=south east] {$3q$};
% Charge at (0,0,0)
\filldraw[blue] (0,0,0) circle (2pt) node[anchor=north west] {$-q$} node[anchor=east, black] {$O$};

% Dashed vertical line for clarity
\draw[dashed] (0,0,0) -- (0,0,1.5) node[midway, right] {$a$};

% Arrangement (b)
% Shifted frame
\begin{scope}[shift={(2,7,1.5)}]
\draw[->] (0,0,0) -- (0,0,2) node[anchor=south]{$z$};
\draw[->] (0,0,0) -- (2,0,0) node[anchor=north east]{$x$};
\draw[->] (0,0,0) -- (0,2,0) node[anchor=west]{$y$};

\node at ((1,-1,2.) {\textbf{(b)}};


% Charge at (0,0,0)
\filldraw[red] (0,0,0) circle (2pt) node[anchor=north west] {$3q$} node[anchor=east, black] {$O$};
% Charge at (0,0,-1.5)
\filldraw[blue] (0,0,-1.5) circle (2pt) node[anchor=north east] {$-q$};

% Dashed vertical line for clarity
\draw[dashed] (0,0,-1.5) -- (0,0,0) node[midway, left] {$a$};
\end{scope}
\end{tikzpicture}
\end{figure}

\noindent
For each arrangement, determine 
\begin{enumerate}
    \item[(i)] The monopole moment, --- \textit{0.2 pts}
    \item[(ii)] The dipole moment, --- \textit{0.4 pts}
    \item[(iii)] The approximate potential (in spherical coordinates) using these moments. --- \textit{0.4 pts}  
\end{enumerate}

\vspace{0.5cm}
\noindent
\underline{Given}: The first two terms of multipole expansion of the potential is given by 
\begin{equation*}
    V(\vec{r}) = \frac{1}{4\pi\epsilon_0} \left( \frac{Q}{r} + \frac{\vec{p}\cdot\hat{r}}{r^2} + \hdots \right),
\end{equation*}
where $Q$ is the monopole moment, $\vec{p}$ is the dipole moment. The spherical coordinates are given by $(r,\theta,\phi)$, where $r$ is the distance from the origin, $\theta$ is the polar angle, and $\phi$ is the azimuthal angle.
 \section*{Solution:}
